% Next-cycle proposal, authored during checkpoint37.
% Complete source: actual_elliptic_full_involution_transporter37.md.
% The checkpoint37 formal slate and all predecessor bytes are unchanged.
\statement{A.133}{Proved}{The exact full elliptic involution transporter and its Brunnian coboundary}
Retain the original negative full-GT involution, its canonical norm
word \(u\), and the actual full profinite elliptic action \(\Theta\).
In the notation of A.131--A.132 put
\[
\begin{gathered}
 a=\sigma^2,\quad b=\tau^2,\quad F_b=\phi(a,b),\quad
 B_u=u(a,b),\\
 U=u(X,K),\quad Z_u=u(X_2,b),\quad
 \Phi=\phi(X,K),\quad Z_\phi=\phi(X_2,b).
\end{gathered}
\]
The actual unique pure transporter of \(\Pi_u\) is
\[
 h_u=F_b,\qquad \Theta(\Pi_u)=F_b\Pi_uF_b^{-1}.
 \tag{A.133.1}
\]
It lies in the entire Brunnian kernel \(\mathcal B\), and its
involution class is explicitly neutral there:
\[
 B_u\in\mathcal B,\qquad
 B_u^{-1}\Theta(B_u)=F_b^{-1}.
 \tag{A.133.2}
\]
The specified conjugate
\[
 \Pi_u^*=B_u\Pi_uB_u^{-1}
        =(UZ_u)\Pi_0(UZ_u)^{-1}
 \tag{A.133.3}
\]
has exact order three and is fixed by the \emph{original}
\(\Theta\). These statements do not assert \(h_u=1\),
the specified elliptic boundary, or the original CM comparison.

\begin{proof}
\emph{1. Full marked inputs and the two cycles of the original parameter.}
All groups and completions are those of A.121 and A.126. Thus
\[
\begin{gathered}
 X_2=\sigma X\sigma,\quad X_3=\tau X_2\tau,\quad
 XX_2X_3=1,\\
 K=[Y,X]=\sigma\tau^2\sigma,\quad K_2=\tau\sigma^2\tau,
 \quad\Pi_0=\tau^{-1}\sigma^{-1}Y.
\end{gathered}
\]
Use \([g,h]=ghg^{-1}h^{-1}\). On the original free pair,
\[
 u=\phi\alpha(u),\qquad
 \alpha(x)=x^{-1},\quad
 \alpha(y)=\phi^{-1}y^{-1}\phi.
 \tag{A.133.4}
\]
Both \(\phi\) and \(u\) satisfy the original pentagon, the latter
by A.37. The original elliptic formulas give
\[
\begin{gathered}
 \Theta(\sigma)=F_b\sigma^{-1}F_b^{-1},\qquad
 \Theta(\tau)=\tau^{-1},\\
 \Theta(X)=\Phi X^{-1}\Phi^{-1},\qquad
 \Theta(K)=K^{-1},\qquad\Theta^2=1.
\end{gathered}                                           \tag{A.133.5}
\]

For clarity, the same original \(\phi\) has both three-cycles
needed below. If \(c\) is simultaneous inversion on the original
pair, direct composition gives
\[
 (\alpha c)(x)=x,\qquad (\alpha c)(y)=\phi^{-1}y\phi.
\]
Thus \(\alpha c=(1,\phi)\) is an actual GT0 element. The
negative and zero three-cycles of this same word, in C.30's
exact convention, are
\[
\begin{gathered}
 \phi(s,t)s^{-1}\phi(d_0,s)d_0^{-1}\phi(t,d_0)t^{-1}=1,\\
 \phi(s,t)\phi(d_0,s)\phi(t,d_0)=1,
 \qquad d_0=(st)^{-1}.
\end{gathered}                                           \tag{A.133.6}
\]
Only \(\phi\) is used in these cycles; no zero or negative
three-cycle for \(u\) is assumed.

\par\smallskip\noindent\emph{2. The exact rank-two braid identity.}
Put
\[
 A_\phi=\phi(X,a),\quad c_a=(Xa)^{-1},\quad
 J_a=\phi(a,c_a).
\]
The zero cycle gives \(\phi(c_a,X)=A_\phi^{-1}J_a^{-1}\).
Its substitution into the negative cycle yields
\[
 A_\phi X^{-1}A_\phi^{-1}
                 =aJ_a^{-1}c_aJ_a.                    \tag{A.133.7}
\]
The relation \((\sigma X)^2=(X\sigma)^2\) makes
\(q_0=XX_2=X\sigma X\sigma\) central in
\(\langle X,\sigma\rangle\). Hence
\(X_2^{-1}=q_0^{-1}X\), while
\(\operatorname{Ad}(\sigma)(c_a)=X_2^{-1}\) and
\(\operatorname{Ad}(\sigma)(a)=a\).
Conjugation equivariance and duality therefore give
\[
 \sigma J_a^{-1}\sigma^{-1}
 =\phi(a,X_2^{-1})^{-1}
 =\phi(X_2^{-1},a)=\phi(X,a)=A_\phi.
\]
The last equality removes a central decoration from a commutator
word; the ordinary exponent-sum argument extends continuously to
its profinite closure. Multiplying (A.133.7) on the left and
right by \(\sigma^{-1}\), and using \(\sigma^{-1}a=\sigma\),
proves
\[
 \sigma^{-1}A_\phi X^{-1}A_\phi^{-1}\sigma^{-1}
       =A_\phi X_2^{-1}A_\phi^{-1}.
 \tag{A.133.8}
\]
This is a literal identity in the full original group.

\par\smallskip\noindent\emph{3. The complete action on the second marked pair.}
Set \(E_\phi=\phi(X_3^{-1},K_2)\). The actual type-B image
of the original pentagon is
\[
                  F_bE_\phi A_\phi=\Phi Z_\phi.
 \tag{A.133.9}
\]
Here \(X\) commutes with \(X_2,b\), so it commutes with
\(Z_\phi\). Also \(\sigma\) commutes with \(X_3,K_2\),
so it commutes with \(E_\phi\). Finally \(b\) commutes with
\(X,K\), hence with \(\Phi\). For the braid assertions use
\[
 Z_c=(\sigma\tau)^3=aK_2=bK,
\]
which is central in \(\langle\sigma,\tau\rangle\), and the
Artin and far-commutation relations for \([\sigma,X_3]=1\).
No global elliptic centrality of \(Z_c\) is claimed.

Let \(H=F_bE_\phi\). Equations (A.133.5) and (A.133.9),
with the preceding centralizations, give
\[
 \Theta(X)=HA_\phi X^{-1}A_\phi^{-1}H^{-1},\qquad
 \Theta(\sigma)=H\sigma^{-1}H^{-1}.
\]
Use \(X_2=\sigma X\sigma\) and (A.133.8). This proves
\[
\begin{aligned}
 \Theta(X_2)
 &=H\sigma^{-1}A_\phi X^{-1}A_\phi^{-1}\sigma^{-1}H^{-1}\\
 &=HA_\phi X_2^{-1}A_\phi^{-1}H^{-1}\\
 &=\Phi Z_\phi X_2^{-1}Z_\phi^{-1}\Phi^{-1}.
\end{aligned}                                           \tag{A.133.10}
\]
Since \([\Phi,b]=1\), the restriction of
\(\operatorname{Ad}(\Phi^{-1})\Theta\) to the actual pair
\((X_2,b)\) sends it exactly to
\((Z_\phi X_2^{-1}Z_\phi^{-1},b^{-1})\).

\par\smallskip\noindent\emph{4. The original norm after actual pair evaluation.}
Equation (A.133.4) has the following word consequence for every
pair \((s,t)\) in any profinite group:
\[
 u\bigl(\phi(s,t)s^{-1}\phi(s,t)^{-1},t^{-1}\bigr)
       =u(s,t)\phi(s,t)^{-1}.
 \tag{A.133.11}
\]
Indeed the substitution is \(\operatorname{Ad}(\phi)\alpha\)
on the abstract original free pair. Conjugating the identity
\(\alpha(u)=\phi^{-1}u\) by \(\phi\) gives the displayed order.
Continuous evaluation then proves the assertion without an
injectivity assumption on the evaluated pair.

Apply (A.133.11) to \((X_2,b)\), using (A.133.10), and to
\((a,b)\), using (A.133.5). Together with its application to
\((X,K)\), this gives
\[
\begin{gathered}
 \Theta(Z_u)=\Phi Z_uZ_\phi^{-1}\Phi^{-1},\qquad
 \Theta(B_u)=B_uF_b^{-1},\\
 \Theta(U)=U\Phi^{-1}.
\end{gathered}                                           \tag{A.133.12}
\]
A.131 and A.132 identify the actual unique pure transporters as
\(p_u=B_u^{-1}UZ_u\) and \(h_1=\Phi Z_\phi\). Thus
\[
\begin{gathered}
 \Theta(UZ_u)=UZ_u(\Phi Z_\phi)^{-1}=UZ_uh_1^{-1},\\
 \Theta(p_u)=F_bp_uh_1^{-1}.
\end{gathered}
 \tag{A.133.13}
\]
Their exact defining formula
\(h_u=\Theta(p_u)h_1p_u^{-1}\) proves (A.133.1).

\par\smallskip\noindent\emph{5. The entire kernel and the two distinct cocycles.}
The based table has
\(q(a)=(1,1,K)\) and \(q(b)=(K,1,1)\).
Every coordinate of \(q(F_b)\) and \(q(B_u)\) therefore has
one commutator-word argument equal to one. Both words lie in
\(P'\), so \(F_b,B_u\in\mathcal B\).
The standard \(C_2\)-cocycle, named distinctly from the cyclic
\(C_3\)-defect, is
\[
 \zeta_{u,2}=h_u^{-1}=F_b^{-1}
                =B_u^{-1}\Theta(B_u).
 \tag{A.133.14}
\]
This is an explicit coboundary inside the entire \(\mathcal B\).
The other sign also has an explicit cobounder:
\(F_b=\Theta(B_u)^{-1}\Theta^2(B_u)\).
No assertion that inversion preserves arbitrary nonabelian
cohomology classes is used. The distinct cyclic defect is
\[
 \zeta_{u,3}=\Theta(\Pi_u)\Pi_u^{-1}=[F_b,\Pi_u]
          \in\mathcal B.
\]
The original element \(h_u\) need not be one.

Put \(r_u=UZ_u\). Equation (A.133.13) proves
\[
 \Theta(r_u)h_1=r_u,\qquad
 \Theta(\Pi_u^*)=\Pi_u^*.
\]
The exact order-three element \(\Pi_u^*\), together with the
original external involution, gives \(C_3\times C_2\cong C_6\)
in \(G\rtimes\langle\Theta\rangle\).
There is also the explicit pure-derived inner gauge
\[
 \operatorname{Ad}(r_u^{-1})\Theta\operatorname{Ad}(r_u)
      =\Theta_0:=\operatorname{Ad}(h_1^{-1})\Theta.
 \tag{A.133.15}
\]
It fixes \(\Pi_0\), hence commutes with
\(\rho_0=\operatorname{Ad}(\Pi_0)\).
Since \([Z_\phi,X]=1\), (A.133.5) and (A.133.10) give
\(\Theta_0(X)=X^{-1}\) and
\(\Theta_0(X_2)=X_2^{-1}\). The same-sign handles commute
and \(XX_2X_3=1\), so also
\(\Theta_0(X_3)=X_3^{-1}\).
This identifies an actual marked action, without identifying
\(\Theta_0\) with a standard elliptic conjugation.

\par\smallskip\noindent\emph{6. Exact remaining comparison scope.}
A.124 proves the full normalizer equality
\(N_G(\langle\Pi_u^*\rangle)=\langle\Pi_u^*\rangle\).
Since \(\Theta\) fixes this element, an element
\(g\Theta^\epsilon\) normalizes its subgroup precisely when
\(g\) does. Hence
\[
 N_{G\rtimes\langle\Theta\rangle}
       (\langle\Pi_u^*\rangle)
   =\langle\Pi_u^*\rangle\times\langle\Theta\rangle.
\]
Every element of this normalizer centralizes the cubic. Thus the
original modular \(S_3\) frame of A.111 cannot map into this
extension with its order-three generator sent to \(\Pi_u^*\)
or its inverse: the original transposition would have to invert it.
This excludes only that specified comparison with the Y-direction
elliptic translation class, not other based sewing constructions.
No original CM conclusion, specified boundary
\([\omega,\xi]=K^d\), or equality \(a(k)=a_B\) follows here.
No full three-strand GT0 action is asserted.
\end{proof}

\paragraph{Source record.}
The complete new full-word proof is retained in
\path{work_A/actual_elliptic_full_involution_transporter37.md}.
The precise normalizer comparison and A.124 proof scope are in
\path{work_A/actual_elliptic_fixed_cubic_normalizer_boundary37.md}.
All original formulas, both original cycles, and every pentagon
pair are retained. Root and DE independently checked the new
ordered calculation before its full saved-source review. No
finite computation, private build, or frozen checkpoint37
formal edit is used by this proposal.
